\chapter{Data Structures}
\label{ch:datastructures}

The previous chapters computed with mathematics; this one steps back to ask what,
exactly, a value \emph{is} inside Mathilda. The answer is unusually simple, and it
is the foundation everything else rests on: there is only one kind of value, the
\emph{expression}, and every list, matrix, rule, pattern, and association you will
ever type is one. That uniformity is what lets a single handful of tools---\B{Part},
\B{Map}, \B{ReplaceAll}---work on everything without special cases. But simplicity at
the surface is not the same as naivety underneath. Behind the uniform tree, Mathilda
keeps two specialized structures that make real work fast: the \emph{association},
a hash-backed key--value store with amortized \(O(1)\) lookup, and the \emph{packed
array}, a dense machine-precision buffer that runs numeric code at the speed of C.
The theme of the chapter is precisely this tension---one uniform data model on top,
invisible acceleration beneath---and how Mathilda keeps the two from ever
contradicting each other.

We build up in four steps. First the expression tree itself: head and arguments,
atoms and compounds, and the tools that read a value's shape. Then lists---the
workhorse aggregate---their parts, spans, and the newer array operations. Then
associations, the key--value store. And finally the packed-array substrate, where
most of the chapter's ``under the hood'' weight lies, because that is where a
transparent list quietly becomes a C buffer and back again.

\section{Everything is an expression}
\label{sec:ds-expressions}

Type \(a + b\,c\) and it looks like an infix formula. It is not stored as one.
Every compound value in Mathilda is a \emph{function application}\index{expression!head
and arguments}---a head applied to a sequence of arguments---and \B{FullForm} prints
that underlying tree with the surface syntax stripped away.

\begin{codepairs}
\pair{06-data-structures/expressions/1}{The sum-of-products \(a+bc\) is
\texttt{Plus} applied to \(a\) and \texttt{Times[b, c]}: operators are just
shorthand for named heads.}
\pair{06-data-structures/expressions/2}{A list is no different---\(\{1,2,3\}\) is
\texttt{List[1, 2, 3]}, the head \B{List} applied to three arguments.}
\pair{06-data-structures/expressions/3}{Even a fraction is a compound: \(1/2\) is
\texttt{Rational[1, 2]}. What prints as a number is a two-argument tree.}
\end{codepairs}

Because there is only one representation, the head of a value tells you what kind of
thing it is, and \B{Head} reads it off for anything at all---an operator expression,
an atom, a list, or a call to an undefined function.

\begin{codepairs}
\pair{06-data-structures/expressions/4}{The head of \(a+bc\) is \B{Plus}\ldots}
\pair{06-data-structures/expressions/5}{\ldots and of the sub-product \(bc\) it is
\B{Times}.}
\pair{06-data-structures/expressions/6}{A machine integer has head \texttt{Integer}\ldots}
\pair{06-data-structures/expressions/7}{\ldots a machine real, head \texttt{Real}\ldots}
\pair{06-data-structures/expressions/8}{\ldots a fraction, head \B{Rational}\ldots}
\pair{06-data-structures/expressions/9}{\ldots and a complex number, head \B{Complex}.}
\pair{06-data-structures/expressions/10}{Strings carry head \texttt{String}.}
\pair{06-data-structures/expressions/11}{A list's head is \texttt{List}\ldots}
\pair{06-data-structures/expressions/12}{\ldots and an unknown \(f[x]\) has head
\(f\), the symbol you named. Nothing is privileged.}
\end{codepairs}

The values that have no internal parts to descend into---numbers, symbols,
strings---are the \emph{atoms}\index{atom}; everything else is a compound with a
head and arguments. \B{AtomQ} draws the line, and it is worth watching where it
falls, because ``looks compound'' and ``is compound'' are not the same.

\begin{codepairs}
\pair{06-data-structures/expressions/13}{An integer is an atom\ldots}
\pair{06-data-structures/expressions/14}{\ldots so is a symbol\ldots}
\pair{06-data-structures/expressions/15}{\ldots and so, despite printing as
\(1/2\), is a \B{Rational}: it is indivisible as a value even though its
\texttt{FullForm} shows two integers.}
\pair{06-data-structures/expressions/16}{A function call is not an atom---you can
take it apart.}
\pair{06-data-structures/expressions/17}{Neither is a list.}
\end{codepairs}

\calloutnote{Under the hood}{An \texttt{Expr} is a tagged union
(\texttt{src/expr.h}): the tag is one of \texttt{EXPR\_INTEGER} (a machine
\texttt{int64}), \texttt{EXPR\_REAL} (a \texttt{double}), \texttt{EXPR\_SYMBOL},
\texttt{EXPR\_STRING}, \texttt{EXPR\_BIGINT} (a GMP integer), or
\texttt{EXPR\_FUNCTION} (a head pointer plus an argument array). The first five are
the atoms; only \texttt{EXPR\_FUNCTION} has parts. Expressions are treated as
immutable during evaluation---a transformation returns a new tree rather than
editing one in place---and \texttt{expr\_copy} is a reference-count bump, not a deep
clone, so this costs far less than it sounds. A \texttt{Rational} and a
\texttt{Complex} are \texttt{EXPR\_FUNCTION} nodes at the level of storage, but
\B{AtomQ} reports them atomic because they behave as single numbers; that deliberate
override is the gap between the two answers above.}

The shape of a tree---how deeply it nests, and what lives at each level---is itself
queryable. \bidx{Depth}\mcode{Depth} gives one more than the longest path from the
root to a leaf, and \B{Level} collects the subexpressions at a chosen depth.

\begin{codepairs}
\pair{06-data-structures/expressions/18}{\mcode{Depth} of \(f[g[h[x]]]\) is
\(4\): three nested heads plus the atom \(x\) at the bottom.}
\pair{06-data-structures/expressions/19}{A \(2\times2\) matrix has depth \(3\)---the
outer list, the rows, and the numbers.}
\pair{06-data-structures/expressions/20}{\B{Level} with the specification
\(\{-1\}\) gathers the atoms (level \(-1\) counts up from the leaves): every
number and symbol at the bottom of \(a + f[x, y^n]\).}
\pair{06-data-structures/expressions/21}{Level \(2\) instead gathers everything
within two steps of the root, so \(y^n\) appears whole rather than split into
\(y\) and \(n\).}
\end{codepairs}

A negative level counts distance from the leaves and a positive one distance from
the root, so \(\{-1\}\) is ``the atoms'' and \(2\) is ``down to two levels
deep.''\index{levelspec} These level specifications are shared by \B{Map},
\B{Cases}, \B{Position}, and the rest of the structural operators---learn them once
and they apply everywhere, which is the recurring dividend of the uniform model.

\section{Lists}
\label{sec:ds-lists}

The list is the aggregate you will reach for most, and because it is just an
\texttt{EXPR\_FUNCTION} with head \texttt{List}, every generic tool applies to it
unchanged. What lists add is a rich vocabulary for addressing their elements.

\subsection{Parts and spans}
\label{ssec:ds-parts}

\B{Part}---written \(\mathit{expr}\)\texttt{[[\ldots]]}---extracts an element by
position, counting from \(1\), or from the end with negative indices. A
\emph{span}\index{list!span}---the \B{Span} operator, written \texttt{i ;; j}---extracts
a contiguous run, optionally with a step; and for a nested list a comma-separated
index descends level by level.

\begin{codepairs}
\pair{06-data-structures/lists/2}{The third element of the vector of primes.}
\pair{06-data-structures/lists/3}{A negative index counts from the end: \(-1\) is
the last.}
\pair{06-data-structures/lists/4}{A span \texttt{2 ;; 4} takes elements two through
four inclusive.}
\pair{06-data-structures/lists/5}{With a step, \texttt{1 ;; -1 ;; 2} takes every
second element from first to last.}
\pair{06-data-structures/lists/7}{On a matrix, \texttt{[[2, 3]]} descends: row
two, then column three.}
\pair{06-data-structures/lists/8}{\texttt{All} in a slot keeps that whole axis, so
\texttt{[[All, 2]]} is the second column.}
\end{codepairs}

The same \texttt{[[\ldots]]} notation reads a single element, a slice, or a whole
axis depending on what you put in the slot, and it composes down through any depth
of nesting---a direct consequence of a matrix being nothing more special than a list
of lists.

\subsection{Shape and nesting}
\label{ssec:ds-shape}

\B{Length} counts the top-level arguments of any expression, and \B{Dimensions}
reports the rectangular shape of a nested list---the lengths at each level, read
outward in.

\begin{codepairs}
\pair{06-data-structures/lists/9}{Six primes, length six.}
\pair{06-data-structures/lists/10}{\B{Length} of the matrix is \(2\)---it has two
rows; length is always the count of immediate arguments.}
\pair{06-data-structures/lists/11}{\B{Dimensions} gives the full shape:
\(2\times3\).}
\pair{06-data-structures/lists/12}{And for a rank-\(3\) array, all three axis
lengths.}
\pair{06-data-structures/lists/13}{Because \B{Length} is generic, it counts the
arguments of any head---here the four arguments of \(f\).}
\end{codepairs}

The contrast between the last two pairs is the point: \B{Dimensions} presumes the
rectangular, tensor-like structure a numeric array has, while \B{Length} makes no
such assumption and works on \(f[a,b,c,d]\) exactly as on a list. One is for data;
both are for expressions.

\subsection{Reshaping, padding, and gradients}
\label{ssec:ds-arrayops}

Four operations turn a flat list into a shaped one, resize it, differentiate it, and
locate a value inside it. \B{ArrayReshape} pours the elements of a list, in
row-major order, into an array of whatever shape you request; \B{ArrayPad} grows (or,
with negative amounts, trims) an array at its edges under a choice of fill scheme;
\B{ListGradient} takes the numerical derivative of sampled data by finite
differences; and \B{FirstPosition} finds where a pattern first matches.

\begin{codepairs}
\pair{06-data-structures/arrayops/1}{\B{ArrayReshape} lays \(1\ldots12\) into a
\(3\times4\) grid, filling row by row.}
\pair{06-data-structures/arrayops/2}{When the data runs short, the third argument
supplies the fill: seven elements into a \(3\times3\) grid, padded with \(x\).}
\pair{06-data-structures/arrayops/3}{\B{ArrayPad} adds two elements of the default
fill (\(0\)) at each end.}
\pair{06-data-structures/arrayops/4}{An asymmetric \(\{1,2\}\) pad with the
\texttt{"Fixed"} scheme repeats the edge values rather than inserting zeros.}
\pair{06-data-structures/arrayops/5}{A negative amount trims: \(-2\) drops two from
each end.}
\end{codepairs}

\B{ListGradient} is the most mathematically interesting of the four---a
from-first-principles port of \texttt{numpy.gradient}. On the interior of the array
it takes a second-order central difference; at each endpoint, where a centered
stencil would run off the edge, it falls back to a one-sided difference.

\begin{codepairs}
\pair{06-data-structures/arrayops/6}{Sample \(x^2\) at \(x=1,\dots,5\)
(\(\{1,4,9,16,25\}\)); the interior values \(4,6,8\) are exactly the derivative
\(2x\) at \(x=2,3,4\), and the ends \(3,9\) are the one-sided approximations.}
\pair{06-data-structures/arrayops/7}{With symbolic input the finite-difference
formula itself is returned---halved central differences inside, full one-sided
differences at the ends.}
\pair{06-data-structures/arrayops/8}{Sampling \(x^3\) and asking for a fourth-order
stencil recovers \(3x^2\) exactly at every point, endpoints included.}
\end{codepairs}

\calloutnote{Theory}{The single kernel behind \B{ListGradient} is Fornberg's
finite-difference-weight recurrence\index{finite differences!Fornberg weights},
which produces the exact weights for a derivative of any order on any (even
non-uniform) grid. A second-order interior stencil is
\((f_{i+1}-f_{i-1})/2h\)---which is why sampling a quadratic gives the derivative
back with no error, since the central difference of \(x^2\) is exact---and the
endpoints drop to first order (numpy's \texttt{edge\_order = 1}). Because the weights
are exact rationals, integer input yields exact \B{Rational} output and symbolic
input a symbolic formula, as the middle pair shows; only a machine-real array takes
the packed floating-point fast path of \S\ref{sec:ds-packed}.}

\B{FirstPosition} reports the position of the first match as a list of indices---the
very indices \B{Part} would use to extract it---searching depth-first through the
whole expression.

\begin{codepairs}
\pair{06-data-structures/arrayops/9}{The value \(7\) sits at position \(\{3\}\).}
\pair{06-data-structures/arrayops/10}{In a nested list the position is a full index
path: \(4\) is at row two, column two.}
\pair{06-data-structures/arrayops/11}{No match returns \(\mcode{Missing["NotFound"]}\)\ldots}
\pair{06-data-structures/arrayops/12}{\ldots unless you supply a default, which is
returned in its place.}
\end{codepairs}

\calloutnote{Pitfall}{\B{FirstPosition} takes its arguments with the attribute
\texttt{HoldRest}, so the default is evaluated only if it is actually
returned---a deliberate design that lets you write an expensive fallback without
paying for it on a hit. The search follows \B{Position}'s conventions exactly
(levels \(\{0,\infty\}\), \texttt{Heads -> True}), which is why a position can be
the empty list \(\{\}\), meaning the whole expression matched.}

\section{Associations}
\label{sec:ds-associations}

A list indexes by position. Often what you have instead is a set of named
fields---a record, a lookup table, a histogram---and forcing that onto positional
storage means remembering that ``price is element three.'' The \B{Association}\index{association!keyed
store}, written \texttt{<|key -> value, \ldots|>}, indexes by \emph{key}: any expression can
be a key, keys are unique, and they keep their insertion order.

\begin{codepairs}
\pair{06-data-structures/associations/1}{An inventory keyed by fruit name. It prints
back in the order written.}
\pair{06-data-structures/associations/2}{Applied like a function,
\(\mathit{assoc}[\mathit{key}]\) looks the key up---the idiomatic accessor.}
\pair{06-data-structures/associations/3}{The \texttt{[[key]]} form does the same
through \B{Part}.}
\pair{06-data-structures/associations/4}{\B{Keys} lists the keys, in order\ldots}
\pair{06-data-structures/associations/5}{\ldots and \B{Values} the corresponding
values.}
\pair{06-data-structures/associations/6}{\B{Lookup} is the explicit accessor.}
\pair{06-data-structures/associations/7}{Its three-argument form supplies a default
for an absent key---the safe way to read data that may be incomplete.}
\pair{06-data-structures/associations/8}{\B{KeyExistsQ} tests for a key's
presence\ldots}
\pair{06-data-structures/associations/9}{\ldots reporting \texttt{False} when it is
missing.}
\pair{06-data-structures/associations/10}{A bare lookup of an absent key returns a
\mcode{Missing} marker naming what was not found, rather than an error.}
\end{codepairs}

Two behaviors set associations apart from a bag of rules. Keys are unique with
last-value-wins, and functions that consume a collection act on the
\emph{values}\index{association!value threading} while keeping keys aligned.

\begin{codepairs}
\pair{06-data-structures/associations/11}{A repeated key collapses: the later value
wins, but the key keeps its first position.}
\pair{06-data-structures/associations/12}{\B{Map} threads over values and returns an
association---keys ride along untouched.}
\pair{06-data-structures/associations/13}{\B{Total} reduces over the values, exactly
as it would over \texttt{Values[inv]}.}
\pair{06-data-structures/associations/14}{\B{AssociationThread} zips a list of keys
against a list of values into an association---the constructor to reach for when the
two arrive separately.}
\end{codepairs}

\calloutnote{Under the hood}{Single-key access is amortized \(O(1)\). A persistent
key\(\rightarrow\)position hash index (\texttt{src/assoc\_index.h}) is cached on the
association and built on the first single-key read, so \B{Lookup}, \B{KeyExistsQ},
\(\mathit{a}[\mathit{key}]\), and \(\mathit{a}[[\mathit{key}]]\) never scan the
entries. The index reuses spare space already on the expression node, so it costs no
extra memory, and it is transparent to \B{SameQ}, \B{Sort}, and \B{Union}, which see
the association exactly as before. The value-threading rule---that \B{Map}, \B{Total},
\B{Select}, \B{Cases}, and their kin operate on values and keep keys aligned---runs
through one shared helper, which is why these operations compose predictably into
data pipelines.}

\mcode{Part} addresses an association by key exactly as it addresses a list by position,
and a part that selects \emph{several} entries returns a \emph{sub-association}---the
chosen entries, keys and all, in their original order.

\begin{codepairs}
\pair{06-data-structures/assoc-part/1}{A small keyed record to slice.}
\pair{06-data-structures/assoc-part/2}{A span selects a contiguous run of entries and
keeps them keyed.}
\pair{06-data-structures/assoc-part/3}{A list of keys picks exactly those, in the order
requested.}
\pair{06-data-structures/assoc-part/4}{Integer positions return the same
sub-association---position and key are two names for one slot.}
\end{codepairs}

\subsection{Associations are atoms}
\label{ssec:ds-assoc-atoms}

An association looks like a container, but to the structural machinery it is a single,
indivisible object---an \emph{atom}. That is the rule that lets one travel safely through
code that knows nothing about associations. Traversal and testing see the \emph{values};
the keys are addressing metadata, reached only through an explicit \B{Key}.

\begin{codepairs}
\pair{06-data-structures/assoc-atoms/1}{A three-entry association.}
\pair{06-data-structures/assoc-atoms/2}{\B{AtomQ} reports \texttt{True}: structurally it
is one node, not a compound of parts to be taken apart.}
\pair{06-data-structures/assoc-atoms/3}{\mcode{Cases} scans the \emph{values}, so a
pattern filter returns matching values, not entries.}
\pair{06-data-structures/assoc-atoms/4}{\mcode{Position} likewise finds a value, and
names where it lives by its key, wrapped in \B{Key}.}
\pair{06-data-structures/assoc-atoms/5}{And \B{ReplaceAll} cannot touch a key: a rule
that would rewrite \texttt{"y"} leaves the association unchanged, because the keys are not
part of the expression it walks.}
\end{codepairs}

\calloutnote{Under the hood}{Atomicity is a deliberate---and recent---design choice.
\mcode{AtomQ}, \mcode{Depth}, \mcode{LeafCount}, \mcode{Level}, \mcode{Cases},
\mcode{Count}, \mcode{MemberQ}, \mcode{FreeQ}, \mcode{Position}, and the level
specifications of \mcode{Map} and \mcode{Apply} all agree to see only an association's
values. The pay-off is safety: because \texttt{/.} can never merge or rewrite keys, and a
pattern scan can never mistake a key for a value, an association carried through a generic
pipeline---as a rule set, a \mcode{Cases} filter, or a user-defined transformation---behaves
predictably and cannot be silently corrupted. When you \emph{do} need to name a key
structurally, \mcode{Position} hands back a \B{Key}$[\,k\,]$ wrapper that \mcode{Part}
accepts directly.}

\subsection{Operator forms, slots, and pipelines}
\label{ssec:ds-assoc-pipeline}

Data work is a pipeline: filter the rows, sort them, pull a field, summarise. Associations
are built for it, through two conveniences. A pure function can read a field by \emph{name}
with a named slot, \texttt{\#field}, and most operators carry a \emph{curried} form that
takes their data last, so stages compose left to right.

\begin{codepairs}
\pair{06-data-structures/assoc-pipeline/2}{\B{Select} keeps the rows whose \texttt{qty}
field---read by the named slot \texttt{\#qty}---exceeds two.}
\pair{06-data-structures/assoc-pipeline/3}{\B{SortBy} orders the same rows by that field.}
\pair{06-data-structures/assoc-pipeline/5}{\mcode{Lookup["pears"]} is an operator form:
it waits for an association, then looks the key up.}
\pair{06-data-structures/assoc-pipeline/6}{Listable heads thread over the values, so
adding a scalar adds it to each value and keeps the keys aligned.}
\pair{06-data-structures/assoc-pipeline/7}{And \mcode{Total} of a \emph{list} of
associations sums them key by key---the reduction a table of records calls for.}
\end{codepairs}

\calloutnote{Under the hood}{A named slot \texttt{\#field} parses as
\B{Slot}\texttt{["field"]}, which on an association reads that key---so
\texttt{\#qty\,>\,2\,\&} is an ordinary pure function that happens to index by name. The
curried forms come from one table-driven mechanism (\texttt{src/opform.c}) that gives some
forty-three heads---\mcode{Lookup}, \mcode{Select}, \mcode{Map}, \mcode{GroupBy},
\mcode{SortBy}, and their kin---an operator form \texttt{f[args][data]} equivalent to
\texttt{f[data, args]}. Together they let a data pipeline be written point-free, each stage
a self-contained operator applied in turn.}

\subsection{Aggregating: Counts, GroupBy, and Merge}
\label{ssec:ds-aggregate}

The reason associations matter for data work is that they are what the aggregation
operators \emph{produce}. \B{Counts} tallies a list into an association from value to
frequency; \B{GroupBy} partitions a list by a key function; and \B{Merge} combines
several associations, reconciling clashing keys with a function of your choosing.

\begin{codepairs}
\pair{06-data-structures/aggregate/1}{\B{Counts} turns a list of observations into
an empirical frequency table---the discrete distribution of the sample.}
\pair{06-data-structures/aggregate/2}{\B{GroupBy} partitions \(1\ldots10\) by
\texttt{EvenQ}: two groups, keyed \texttt{False} and \texttt{True}.}
\pair{06-data-structures/aggregate/3}{Its three-argument form applies a reducer to
each group---here \B{Total}, collapsing the split-apply-combine in one call.}
\pair{06-data-structures/aggregate/5}{With a \(\mathit{key}\to\mathit{value}\) rule,
\B{GroupBy} groups the rows by their \B{First} field but collects the \B{Last}: sales
gathered per category.}
\pair{06-data-structures/aggregate/6}{Add a reducer and the same call groups,
extracts, and totals---a one-line group-and-sum.}
\pair{06-data-structures/aggregate/7}{\B{Merge} unions the keys of two associations
and applies \B{Total} to the values collected under each: the shared key \(y\) sums
its two contributions, the unshared keys pass through.}
\end{codepairs}

The middle pairs are the split-apply-combine idiom\index{split-apply-combine} that
data analysis lives on: \emph{split} the data into groups, \emph{apply} a summary to
each, \emph{combine} the summaries into one table. \B{GroupBy} does all three, and
because the result is an association it feeds straight into another aggregation---the
frequency table from \B{Counts} is exactly the empirical distribution the statistics
chapter (\S\ref{sec:statistics}) builds its estimators on.

\B{Merge} is one of a family that treats an association as a keyed \emph{set}. The keys
can be intersected, subtracted, and unioned; entries can be dropped by a predicate on
their values; and two tables can be joined on a shared field the way a database would.

\begin{codepairs}
\pair{06-data-structures/assoc-keyalgebra/1}{\B{KeyIntersection} restricts each
association to the keys they all share---here only \texttt{"b"}, each keeping its own
value.}
\pair{06-data-structures/assoc-keyalgebra/2}{\B{KeyComplement} keeps the first
association's keys that appear in no other.}
\pair{06-data-structures/assoc-keyalgebra/3}{\B{KeyUnion} aligns several associations to
their combined key set, filling a gap with a \mcode{Missing} marker.}
\pair{06-data-structures/assoc-keyalgebra/4}{\B{Discard} is \mcode{Select}'s complement:
it drops the entries whose value satisfies the test.}
\pair{06-data-structures/assoc-keyalgebra/5}{\B{JoinAcross} joins two lists of records on
a shared key, merging each matched pair into one row---a relational join.}
\pair{06-data-structures/assoc-keyalgebra/6}{\B{PositionIndex} inverts a list into an
association from each value to the positions it occupies.}
\pair{06-data-structures/assoc-keyalgebra/7}{And \mcode{Merge} accepts a plain list of
rules, collecting clashing keys before it reduces---so it doubles as a grouping tally.}
\end{codepairs}

\usagebox{Association}

\section{NDArrays and packed arrays}
\label{sec:ds-packed}

Everything so far treated a list as a tree of independent expression nodes. For
symbolic work that is exactly right. For numeric work it is ruinous: a list of a
million reals is a million separate nodes, and the evaluator sweeps every argument of
every call on every pass, so merely \emph{holding} a large list costs time on each
operation, whether or not anything is computed. Mathilda's answer is the
\emph{packed array}\index{packed array}---the same list, stored as one dense
machine-precision buffer instead of a million nodes---and the discipline that keeps
that storage change completely invisible.

\subsection{The invisible buffer}
\label{ssec:ds-invisible}

\B{ToNDArray} packs a list into a buffer and hands it back looking exactly like the
list it was. The only tool that can tell the difference is \B{NDArrayQ}, which asks
about storage rather than about value.

\begin{codepairs}
\pair{06-data-structures/packing/1}{\B{ToNDArray} returns what looks like an ordinary
list of reals---because that is what it still is.}
\pair{06-data-structures/packing/2}{Its head is \texttt{List}, it is a list, and it
is not an atom: every structural predicate answers as it would for the unpacked
list.}
\pair{06-data-structures/packing/3}{Only \B{NDArrayQ} reveals that it is packed.}
\pair{06-data-structures/packing/4}{And it is \B{SameQ} to the plain list---same
value, same identity.}
\pair{06-data-structures/packing/5}{\B{DataType} reports the element type\index{data
type!machine (dtype)} of the buffer: a \texttt{"float64"} of doubles.}
\end{codepairs}

The contract is exact: packing changes representation and \emph{nothing else}---not
the value, not an element's head, not the printed form, not the ordering. So the
representation Mathilda picks on its own can never be observed. That single rule is
what the rest of this section is really about.

\subsection{The exactness contract}
\label{ssec:ds-exact}

The sharpest test of the contract\index{packed array!exactness contract} is exact
integers. A list of integers packs to an \texttt{int64} buffer---and every element
must still come back an \emph{Integer}, not a real, because \(1 \ne 1.0\) is a
difference the buffer is forbidden to introduce.

\begin{codepairs}
\pair{06-data-structures/packing/7}{An all-integer list packs to \texttt{"int64"},
not to floats.}
\pair{06-data-structures/packing/8}{And an element read back out has head
\texttt{Integer}, exactly as from the unpacked list.}
\pair{06-data-structures/packing/9}{A reduction stays exact: the sum of \(1\ldots
1000\) is the integer \(500500\), computed on the buffer but with no loss of
exactness.}
\pair{06-data-structures/packing/10}{Where a buffer cannot hold the exact answer, it
abandons: \(\mathit{Range}[10]/2\) is a list of exact \B{Rational}s, which no
machine buffer represents, so the ordinary list gives them.}
\pair{06-data-structures/packing/11}{And a list holding a \B{Rational} refuses to
pack at all---\B{ToNDArray} returns it unchanged rather than round it.}
\end{codepairs}

The rule is the same one that governs the numeric tower everywhere in Mathilda:
compute the buffer answer only when it is \emph{the} answer the interpreter would
give, and otherwise fall back to the slower, exact path. An overflow past
\texttt{int64}, a rational, a radical, a symbol---each makes the buffer decline and
the ordinary list answer.

\calloutnote{Pitfall}{Only \B{ToNDArray} widens a mixed integer/real list to a float
buffer, because you explicitly asked for a single machine type. \emph{Automatic}
packing never does: it leaves \(\{1, 2, 3.\}\) an ordinary list, since silently
turning the exact \(1\) into \(1.0\) would be an observable change and the system is
forbidden to make one on its own. The asymmetry is the contract in action---a
representation you request may cost exactness; one the system chooses for you never
may.}

\subsection{When packing happens}
\label{ssec:ds-auto}

You rarely call \B{ToNDArray} yourself. The common case is that a \emph{producer}
builds a packed array directly: \B{Range}, \B{ConstantArray}, \B{Table} with a
machine-real body, and a handful of others write the buffer without ever
materializing the nodes---which is where most of the win is, since building a million
nodes and then packing them costs far more than writing a million doubles.

\begin{codepairs}
\pair{06-data-structures/packing/12}{\B{Range} of reals packs automatically once the
result is big enough.}
\pair{06-data-structures/packing/13}{A list written out literally in the source never
packs, at any size---there is no producer to opt in, and a source literal is left
exactly as typed.}
\pair{06-data-structures/packing/14}{The master switch \B{\$AutoArrayPacking} turns
automatic packing off\ldots}
\pair{06-data-structures/packing/15}{\ldots and back on. It never touches an explicit
\B{ToNDArray}---that is the user asking, not the system guessing.}
\end{codepairs}

The distinction in the middle pairs surprises people: \mcode{Range[1., 300.]} packs
but the identical-looking literal \(\{1., 2., 3., 4.\}\) does not. Packing is opt-in
per producer, and a hand-written literal has no producer to opt in. A threshold of
four elements (the size of a \(2\times2\) matrix, so that any matrix packs) keeps the
mechanism from paying its own overhead on tiny lists.

\subsection{The visible NDArray}
\label{ssec:ds-visible}

The packed list is one of two faces of the same storage. The other is the
\emph{visible} \B{NDArray}---the identical dense buffer, but wearing its own head, so
it announces itself as a distinct, purely numeric object modelled on numpy's
\texttt{ndarray}. The two share every byte of representation and differ only in what
they claim to be.

\begin{codepairs}
\pair{06-data-structures/ndarray/2}{\B{NDArray} prints wrapped in its own head\ldots}
\pair{06-data-structures/ndarray/3}{\ldots whose \B{Head} is \texttt{NDArray}, not
\texttt{List}.}
\pair{06-data-structures/ndarray/4}{So it is \emph{not} \B{SameQ} to the plain
list---the head is part of a value's identity.}
\pair{06-data-structures/ndarray/5}{\B{NDArrayQ} is \texttt{True} for both faces: it
asks about storage.}
\pair{06-data-structures/ndarray/6}{\B{PackedArrayQ} is the narrower test---\texttt{True}
only for the packed \emph{list}, \texttt{False} for the visible \texttt{NDArray},
which is an atom rather than a list that happens to be packed.}
\pair{06-data-structures/ndarray/7}{\B{Normal} converts a visible array back to an
ordinary nested list.}
\end{codepairs}

The difference is one of intent. A packed list is Mathilda saying ``this is an
ordinary list; I have merely chosen to store it efficiently.'' A visible
\B{NDArray} is you saying ``this is a numeric array; treat it as one''---so it
combines with scalars by numpy-style broadcasting, refuses symbolic operands, and
stays a closed system under the linear-algebra routines. \B{ToPackedArray} and
\B{FromPackedArray} are Mathematica's names for \B{ToNDArray} and \B{FromNDArray},
the same builtins under a second registration.

\begin{codepairs}
\pair{06-data-structures/ndarray/8}{\B{ToPackedArray} widens a mixed integer/real
list to a single \texttt{"float64"} buffer\ldots}
\pair{06-data-structures/ndarray/9}{\ldots the integers becoming doubles, since one
buffer holds one machine type.}
\pair{06-data-structures/ndarray/11}{A derived array stays packed: adding a scalar to
a packed list gives a packed list.}
\pair{06-data-structures/ndarray/12}{So does mapping an elementary function over
it---\B{Sin} of a packed buffer is a packed buffer.}
\pair{06-data-structures/ndarray/13}{And \B{N} widens a packed \texttt{int64} buffer
to a packed \texttt{float64} one in a single pass, rather than boxing a million
reals.}
\pair{06-data-structures/ndarray/15}{The visible \B{NDArray} carries its rank and
shape directly, so \B{Dimensions} reads them in \(O(1)\).}
\end{codepairs}

\usagebox{NDArray}

\subsection{The transparency gate}
\label{ssec:ds-gate}

If a packed list can flow into any operation, what stops an operation that knows
nothing about buffers---a pattern match, a \B{Cases} scan, a user-defined rule---from
reading garbage out of one? The answer is the \emph{transparency
gate}\index{packed array!transparency gate}, and it is the mechanism that makes the
whole abstraction safe.

\begin{codepairs}
\pair{06-data-structures/gate/2}{Start with a packed range.}
\pair{06-data-structures/gate/3}{\B{Cases}---which has no idea packing exists---filters
it with a pattern and gets exactly the right answer.}
\pair{06-data-structures/gate/4}{\B{Reverse} understands buffers, so it reverses in
place\ldots}
\pair{06-data-structures/gate/5}{\ldots and keeps the result packed.}
\pair{06-data-structures/gate/6}{And the buffer path and the ordinary path agree on
the answer, by construction.}
\end{codepairs}

\calloutnote{Under the hood}{The gate lives in the evaluator (\texttt{src/eval.c}),
just before a call reaches its builtin. Its rule inverts the usual default: a packed
buffer is \emph{materialized} back into ordinary nodes for any head that has not
explicitly declared itself packed-aware on the \texttt{AWARE} list in
\texttt{src/pack.c}. So \B{Cases}, \B{Position}, \B{ReplaceAll}, pattern matching,
and every user rule are correct with no code of their own---they simply never see a
buffer. A head like \B{Reverse}, \B{Total}, \B{Sin}, or \B{Dot} that \emph{has} opted
in and been checked receives the buffer intact and runs its C kernel. Forgetting to
opt a head in costs a materialization---a slowdown---never a wrong answer. The
storage overlay is arranged so that a walker which slipped past the gate and iterated
a buffer as if it were a node array would segfault immediately rather than return
something plausible, turning the dangerous silent-wrong-answer class into a loud
crash.}

\calloutnote{Performance}{The payoff is large and the design record measures it on a
\(10^6\)-element list of reals: \B{Length} falls from about \(30.6\)
milliseconds on an ordinary list to \(2\) microseconds packed, \B{Total} from
\(95.8\) milliseconds to \(4.0\), and \B{Sin} from \(320\) milliseconds to
\(16.8\)---the \(30\)-nanosecond-per-element-per-pass evaluator sweep replaced by a
tight loop over contiguous memory (\texttt{docs/design/packed\_arrays.md}). Because a
packed buffer is exactly the representation \B{Compile} wants, these same arrays cross
the compiled-code boundary without being unpacked and repacked at each end---the
subject of Chapter~\ref{ch:compilation}.}

\section*{Where this connects}

Data structures are the substrate the rest of Mathilda stands on. The uniform
expression tree of \S\ref{sec:ds-expressions} is what makes the pattern matcher and
rule engine of Chapter~\ref{ch:programming} possible: a pattern is an expression, a
rule is an expression, and matching one tree against another is the single mechanism
behind \B{Map}, \B{Cases}, \B{ReplaceAll}, and every definition you write. The
levelspecs and structural operators introduced here reappear there as the working
vocabulary of functional programming. Associations are the natural output of the
grouping and counting that descriptive statistics (\S\ref{sec:statistics}) runs on,
and their \(O(1)\) keyed store underlies any code that indexes data by name. The
packed-array substrate of \S\ref{sec:ds-packed} is where numeric Mathilda gets its
speed: it feeds the machine-precision linear algebra of \S\ref{sec:linalg}, and it is
half of the story of \B{Compile} (Chapter~\ref{ch:compilation}), the other half being
the bytecode VM that consumes those buffers. For the exhaustive options of any
function named here, its reference page is one \texttt{?Name} away at the prompt---and
the full machinery of the evaluator, the matcher, and the transparency gate is laid
out in Chapter~\ref{ch:internals}.
